% =====================================================================
% Related Literature Review, v6
% DRAFT (Claude, 2026-11-10, with Anders). Built on the v4 TEXT (Anders:
% "the v4 text was good, I just wanted it more like Bore et al."). v4's
% sentences are kept wherever they are correct. Changes: the six-part
% structure agreed on 2026-11-09 with numbered subsections and run-in
% headings as in Bore and Brensholm (2026), factual errors in v4 corrected
% against the paper notes, a short part on price volatility and persistence,
% and about ten numbers in total (Bore and Brensholm use about ten). See
% drafts/literature-review-v6/CHANGES.md. Text for the students to rework.
% Citations: natbib commands through biblatex natbib=true; keys from
% ../pumped-hydro-thesis/references.bib.
% =====================================================================

\section{Related Literature Review}
\label{sec:literature}

The literature on pumped-storage hydropower addresses several related
questions: whose value is assessed and which capacities are varied, how
uncertainty about future conditions is represented, which markets
provide revenue, and how the resulting planning problems are solved.
This chapter reviews these questions in separate thematic sections, in
which the studies are cited, compared, contrasted and critiqued, and
connected to the present study. Section~\ref{sec:lit_scope} defines the
scope of the review. Section~\ref{sec:lit_valuation} examines valuation
perspectives and plant capacities, Section~\ref{sec:lit_uncertainty} the
representation of uncertainty and price behaviour, and
Section~\ref{sec:lit_markets} market participation.
Section~\ref{sec:lit_methods} then reviews the solution methods applied
to these problems and the physical detail they allow. Finally,
Section~\ref{sec:lit_gap} draws the findings together, identifies the
research gap and positions the present study.

% ---------------------------------------------------------------------
\subsection{Scope and Selection of Literature}
\label{sec:lit_scope}

The review focuses on operational scheduling and economic valuation of
pumped-storage hydropower. Particular attention is given to studies that
connect short-term operating opportunities with longer-term reservoir
management. Studies of conventional reservoir hydropower are included
where they offer relevant insights into water values, uncertainty or
reserve markets, since a reservoir plant faces the same storage decision
without the pumping option. System-planning and hybrid-energy studies
provide context for evaluating storage under different objectives, and
studies of electricity-price behaviour are included because the value
of storage depends on how prices develop over time.

\todonote{TODO (authors): Describe the search process you actually used:
databases (e.g.\ Google Scholar, Scopus), search terms, inclusion
criteria (peer-reviewed, relevance of title and abstract), and whether
further studies were found by following reference lists (backward
snowballing). Three to five sentences are enough. The search log
(\texttt{literature/search-log.md}) is still empty.}

The chapter provides a narrative synthesis of the selected studies. Each
section addresses one dimension of the literature and compares the
studies along that dimension only. Individual studies may therefore
appear in several sections, but each discussion concerns the aspect
relevant to that section. Table~\ref{tab:lit_comparison} in
Section~\ref{sec:lit_gap} summarises the studies across all dimensions.

% ---------------------------------------------------------------------
\subsection{Valuation of Pumped Storage and Its Capacities}
\label{sec:lit_valuation}

The perspective adopted in a study determines what is meant by the value
of pumped storage, and the plant characteristics it varies determine
which part of that value it can explain. This section first examines
valuation perspectives and then the capacities that the studies allow
to change.

\subsubsection{Valuation Perspective}
\label{sec:lit_perspective}

System-oriented studies assess the contribution of storage to the wider
electricity system, whereas owner-oriented studies assess the economic
benefit available to an individual producer. These objectives need to be
distinguished when comparing conclusions about the attractiveness of
storage.

System-oriented studies relate storage value to the surrounding
electricity system. \citet{helseth2013} examine pumped storage within a
hydrothermal system that also includes wind generation and transmission
constraints. \citet{helseth2017} assess an expansion of the
\SI{200}{\mega\watt} Duge pumped-storage plant by up to
\SI{2000}{\mega\watt} through changes in expected socio-economic surplus.
Both studies consider effects beyond the income of an individual plant.
A favourable system-level assessment therefore does not establish the
corresponding private return, since benefits can also accrue to
consumers and other producers. The two can even diverge: the price
differences that pumped storage trades on shrink as more storage is
built \citep{dallinger2019socioeconomic}, and \citet{henden2016economic}
find that a large expansion of pumped storage in Norway leaves its
owners with little income.

Owner-oriented studies instead evaluate the income associated with
operating a particular asset. \citet{ak2019quantifying} report revenue
gains from operating an existing hydropower cascade in pumped-storage
mode, with the gains varying between the investigated price years.
Similarly, \citet{alic2023optimal} assess the economic benefit of adding
a \SI{38}{\mega\watt} pump to a Norwegian cascade. These studies address
the producer's economic opportunities, although their findings remain
specific to the assets and conditions investigated. Increased operating
revenue alone also does not establish investment profitability, which
requires the associated costs to be considered.

Within an owner perspective, the treatment of market power introduces a
further distinction. \citet{abgottspon2013stochastic} consider a
price-making producer whose trading decisions can influence market
outcomes. In contrast, \citet{helseth2016} consider a price-taking
producer facing externally determined prices. This difference affects
the interpretation of flexibility: a price taker responds to available
price differences, while a price maker also considers the price effects
of its own actions.

An owner-oriented valuation can also extend from operating income to an
investment assessment. \citet{bore2026assessing} use additional
hydropower revenues to estimate the maximum investment cost that a
capacity upgrade can support. This distinguishes an assessment of
operating benefits from a valuation that also specifies an economic
lifetime, discount rate and operating-cost assumptions. System benefits,
operating income and investment profitability are thus different
measures of value, and a conclusion about the economic value of storage
must be read against the perspective of the study it is drawn from.

\subsubsection{Pumping, Generating and Reservoir Capacities}
\label{sec:lit_capacities}

Capacity studies differ in which plant characteristics are fixed and
which are allowed to change. Optimising the operation of an existing
plant, comparing specified upgrades and optimising the plant's
dimensions address different questions. Distinguishing these tasks is
essential when interpreting conclusions about the value of additional
capacity.

Studies of specified assets primarily establish how the available
flexibility can be used. \citet{naval2025optimization} optimise the
operation of a given hybrid-energy configuration, whereas
\citet{ak2019quantifying} compare an existing cascade with operation in
pumped-storage mode. \citet{alic2023optimal} and
\citet{schuller2024practical} similarly evaluate a single pump added to
a Norwegian watercourse. These comparisons demonstrate the value of
introducing a pumping option under the investigated conditions. They do
not establish how that value changes when pumping power, generating
power and reservoir capacity are varied independently or jointly.

Comparisons of expansion alternatives address capacity changes more
directly. \citet{helseth2017} assess transmission and pumped-storage
expansions, and find that a detailed representation of the watercourse
lowers the marginal value of additional pumped-storage capacity. The
paper does not state whether pump and turbine capacity are expanded
separately. \citet{helseth2017iet} compare two generator configurations
for a hydropower producer that sells reserve capacity. Both evaluate
changes in available flexibility, but the characteristics varied and the
valuation settings differ. These studies provide precedents for
comparing alternatives without establishing a common ranking of pump,
turbine and reservoir capacities for pumped-storage plants.

\citet{bore2026assessing} examine generating-capacity upgrades of 30, 60
and \SI{90}{\mega\watt} in a conventional hydropower case. They find
diminishing incremental value per unit of additional capacity, with
reserve-capacity revenues contributing most of the gains. This is a
close precedent for assessing capacity changes, but the alternatives
concern generating-power upgrades. They do not compare combinations of
pumping power, generating power and reservoir size in a pumped-storage
plant.

Design optimisation brings selected dimensions into the decision problem
itself. \citet{zisos2023mixing} optimise reservoir dimensions and
photovoltaic capacity in a hybrid island-energy system while holding
pumping and generating power fixed at \SI{6}{\mega\watt}. Their
optimised reservoir is much smaller than the preliminary design, and
uncertainty produces a range of storage sizes and economic outcomes.
This shows that storage sizing can matter for economic performance.
However, the plant is operated by a fixed rule without electricity
prices, and fixing pump and turbine power leaves their joint sizing with
the reservoir outside the examined design choices.

The reviewed studies therefore provide evidence on operating given
assets, introducing pumping, comparing expansions and selecting
reservoir dimensions. They differ in which capacities are varied and
which remain fixed. This distinction matters because an additional unit
of one capacity may have limited value when another capacity is
binding. A larger pump, for example, offers limited additional storage
opportunity when the reservoir is already full. None of the reviewed
studies varies pumping power, generating power and reservoir capacity
jointly for a plant that trades in electricity markets.

% ---------------------------------------------------------------------
\subsection{Representation of Uncertainty}
\label{sec:lit_uncertainty}

The value of stored water depends on prices and inflows that are unknown
when operating decisions are made. The reviewed studies differ in which
inputs they treat as uncertain, how these inputs are described, and when
the uncertainty is assumed to be resolved. All three choices affect what
a reported storage value represents. In addition, the value of storage
depends on how prices behave over time, which few of the studies
examine.

\subsubsection{Sources of Uncertainty}

The sources of uncertainty follow the question a study addresses.
Owner-oriented studies concentrate on inflows and energy prices
\citep{alic2023optimal,bore2026assessing}, and \citet{helseth2016}
additionally treat reserve-capacity prices as uncertain. System and
hybrid-energy studies emphasise the renewable resource instead:
\citet{helseth2013} include wind power alongside inflow,
\citet{naval2025optimization} consider wind generation, and
\citet{zisos2023mixing} consider wind, demand and the wind-to-power
conversion. \citet{helseth2017} combine weather uncertainty with
exogenous price uncertainty. Results obtained under one set of uncertain
inputs therefore cannot simply be transferred to another. For a
price-taking owner, price uncertainty is the most direct driver of the
value of shifting energy in time, yet it receives no attention in the
hybrid-system studies, where prices play a subordinate role.

\subsubsection{Describing Uncertainty over Time}

The description of the uncertain inputs also varies.
\citet{ak2019quantifying} evaluate annual revenues for historical inflows
and a set of electricity-price years. This reveals how revenue gains
vary between price years, but each year is evaluated on its own, so the
results do not show how an operator should act without knowing which
year will occur. \citet{naval2025optimization} use chance constraints,
which require operating targets to be met with a specified probability.
Such a formulation addresses the reliability of a schedule under wind
uncertainty; it does not describe how decisions adapt as outcomes are
observed. Other studies describe uncertainty as a process that develops
between decision stages. \citet{helseth2017iet} and
\citet{alic2023optimal} use weekly Markov states, and
\citet{bore2026assessing} construct a Markov chain from clustered
historical data. A Markov description retains the dependence between
consecutive weeks. This property matters for storage decisions, because
the benefit of holding water depends on how long high- or low-price
conditions are expected to last.

The dependence between the uncertain inputs is a further choice.
\citet{kleiven2023} find that ignoring the negative correlation between
Norwegian prices and inflow lowers water values and changes reservoir
operation, but has little effect on expected revenue.
\citet{perezdiaz2020medium} find the correlation between energy and
reserve prices too weak to affect water values for a Spanish plant.
Such findings depend on the data and the plant, however, and say little
about whether the same holds when prices also drive pumping.

A further distinction concerns when information becomes available to the
decision maker. In the weekly-stage studies of \citet{alic2023optimal},
\citet{helseth2016} and \citet{helseth2017iet}, new information arrives
every week and later decisions respond to it. \citet{helseth2017} also
revisit decisions every week, but treat the uncertainty beyond the
current week as resolved in a single step. The two descriptions imply
different degrees of foresight about the more distant future, even when
the evaluation period is the same. The length of the horizon alone
therefore says little about how realistically a study represents
decision-making under uncertainty.

\subsubsection{Information within the Week}

Uncertainty within a decision stage is a separate issue from uncertainty
between stages. \citet{bore2026assessing} assume that prices and inflow
for the current week are known when its decisions are made. Their
framework thus combines uncertainty between weeks with perfect foresight
within a week, which allows short-term operation to be timed better than
is possible in practice. \citet{abgottspon2014} examine this
simplification directly by comparing known and uncertain descriptions of
operation within the week for a Swiss pumped-storage plant. Treating the
week as known overvalues water, but describing it as uncertain improves
expected profit only when weekly secondary reserves are included. This
suggests that within-week foresight is a less restrictive assumption for
energy trading than it first appears, but not when reserves are sold.
The conclusion rests on a single plant whose capacities are not
reported.

\subsubsection{Price Volatility and Persistence}

A statement that a study ``includes uncertainty'' is therefore not
sufficient to characterise it. For a storage owner trading on price
differences, the representation of price behaviour over time is likely
to be at least as important as the number of uncertain inputs. Two
properties of this behaviour are central: volatility, how much prices
vary around their expected pattern, and persistence, how long high- or
low-price conditions last (Section~\ref{sec:ph25-volatility-persistence}).
\citet{haldrup2006regime} find that Nordic day-ahead prices are highly
persistent, most of all in the Norwegian price areas, which they link to
the slow variation of reservoirs and inflow. A weekly Markov description
captures such persistence only through its transition probabilities, and
\citet{mo2001} note that it therefore forgets deviations from the
expected price level quickly.

Studies that relate price behaviour to storage value do so with known
prices. \citet{ak2019quantifying} find the largest gain from pumping in
the price year with the largest variation within the day, although the
relation is not one-to-one. \citet{sioshansi2009estimating} find that
most of the arbitrage value in the US PJM market is earned by the first
hours of storage, and \citet{bradbury2014economic} find that higher
volatility raises the return on storage but is not the main driver of
the most profitable ratio of storage to power.
\citet{lohndorf2023} report that coordinating day-ahead and intraday
trading is worth about an order of magnitude more for a pumped-storage
plant than for a reservoir plant without a pump. None of these studies
varies volatility and persistence separately or relates them to the
value of pump, turbine and reservoir capacity under uncertainty. Since
pumping and generating power earn from short price differences, whereas
reservoir capacity earns from differences that last over weeks or
seasons, the two properties may favour different capacities.

% ---------------------------------------------------------------------
\subsection{Electricity Market Participation}
\label{sec:lit_markets}

Market participation determines the revenue opportunities available to
a storage owner. Energy-market participation rewards electricity sales
relative to the cost of purchases for pumping. Other products can
remunerate the availability of operating flexibility or its actual use.
Studies must therefore be compared by the products and settlements
included, since these define what their reported revenues cover. The
reported gains are not directly comparable either, as market design,
price levels and plant characteristics differ between the studies.

\paragraph{Energy markets}
Energy-centred studies provide a baseline for assessing the value of
pumping. \citet{ak2019quantifying} examine revenue gains associated with
electricity-price variations and find that pumped-storage operation
raises net revenue by about 6\% on average, while \citet{alic2023optimal}
evaluate price-based operation under environmental constraints. Their
focus differs from studies that include additional trading
opportunities. For example, \citet{abgottspon2013stochastic} consider
both day-ahead and forward trading. Forward positions extend the
producer's commercial opportunities, but they serve a different purpose
from payments for maintaining reserve availability.

\paragraph{Reserve capacity}
Reserve-capacity participation creates an additional use for the same
operating flexibility. \citet{helseth2016} coordinate sales into
separate, sequentially cleared energy and capacity markets and examine
the implications for schedules and water values.
\citet{abgottspon2014} also include ancillary services in
pumped-storage scheduling. In related conventional-hydropower work,
\citet{helseth2017iet} find that reserve sales change operation and
increase operating profit by 7 to 11\%, and \citet{perezdiaz2020medium}
report a gain of about 4.5\% for a Spanish plant selling automatic
frequency restoration reserves. These studies show that reserve
participation can affect the allocation of resources as well as provide
an additional source of revenue. For pumped storage, the pump widens
this flexibility: \citet{dallinger2019socioeconomic} show that
pump-turbines let a plant offer downward reserve without generating, so
its reserve provision no longer depends on natural inflow.

\paragraph{Reserve activation}
The treatment of activation further distinguishes market
representations. \citet{chazarra2017} include regulation energy
associated with upward and downward reserve use, alongside energy and
reserve scheduling, and find that reserve capacity is the main source of
income for a closed-loop plant. Accounting for regulation energy
connects reserve provision to actual electricity consumption or
generation. By comparison, an assessment limited to availability
payments does not capture all the associated energy settlements and
changes in stored water.

Broader market coverage is considered by \citet{bore2026assessing}, who
include day-ahead and intraday energy alongside FCR-N, FCR-D and upward
and downward aFRR capacity. Actual reserve activation is excluded. Their
results indicate that capacity-market participation affects reservoir
management and strengthens revenues in periods with low energy prices.
Compared with the regulation-energy treatment in \citet{chazarra2017},
this example highlights that the number of included products and the
completeness of their delivery representation are separate dimensions
of market coverage.

Consequently, studies of market participation need to distinguish energy
trading, reserve-capacity commitments and reserve activation. Revenues
from these activities interact because they depend on shared operating
capability. Findings from different market combinations cannot be
compared solely through their total income, particularly when their
treatment of activation and settlement differs. Most of the multi-market
evidence also concerns conventional reservoir hydropower. For pumped
storage, where the pump offers an additional means of providing
reserves, the long-term effect of reserve participation is less well
documented.

% ---------------------------------------------------------------------
\subsection{Solution Methods}
\label{sec:lit_methods}

The preceding sections show that the valuation of pumped storage can
combine a long planning horizon, uncertain inputs, physical detail and
several markets. No solution method accommodates all of these to the
same degree, and the physical detail a study can retain depends on the
method it uses. This section reviews the methods applied in the
literature, whose principles are introduced in
Sections~\ref{subsec:two_stage_stochastic_optimisation},
\ref{subsec:stochastic_dynamic_programming} and
\ref{subsec:stochastic_dual_dynamic_programming}, and compares what each
permits and requires.

\subsubsection{Deterministic and Scenario-Based Optimisation}

The most direct approach formulates the scheduling problem as a single
optimisation problem with given inputs. \citet{ak2019quantifying} use
nonlinear programming together with monthly reservoir operating rules,
and solve the problem separately for each price year.
\citet{chazarra2017} formulate a mixed-integer model for short-term
scheduling, which allows the discrete operating modes of the plant to be
represented. This includes hydraulic short-circuit operation, in which
pumping and generation take place simultaneously, and which they find
can greatly increase the income of the plant.
\citet{naval2025optimization} extend a mixed-integer model with chance
constraints to account for wind uncertainty in hourly operation over a
year.

These formulations place few restrictions on the plant description,
which explains their use where nonlinear or discrete characteristics are
central. Their limitation lies in the treatment of time and information.
The schedule is optimised for the inputs supplied, and the long-term use
of the reservoir is governed by prescribed rules or targets, as in
\citet{ak2019quantifying}. The methods are therefore well suited to
short horizons or to evaluating given conditions, but they do not by
themselves provide the value of stored water under uncertainty.

\subsubsection{Two-Stage Stochastic Programming and Rolling Horizons}

Two-stage stochastic programming introduces uncertainty while keeping a
single optimisation problem. \citet{helseth2017} embed two-stage
stochastic linear programmes in a rolling-horizon simulator: each weekly
decision is optimised against a set of scenarios for the future, and the
problem is solved again as the simulation advances. The authors choose
this approach to handle many reservoirs without aggregation and to use
historical data directly. Compared with the deterministic formulations
above, the weekly decisions are informed by a range of possible futures.

The approach has two limitations. First, the scenarios are treated as
fully known after the first-stage decision, so the future part of each
problem is operated with foresight. Second, the computational effort
grows with the number and length of the scenarios.
\citet{helseth2017} accordingly reduce the number of scenarios and use a
coarser time resolution along them to limit computation time, and still
report run times of several days. A rolling horizon thus offers
flexibility in the system description at the expense of a coarser
representation of how information is revealed.

\subsubsection{Stochastic Dynamic Programming}

Stochastic dynamic programming (SDP) divides the problem into stages and
computes the expected future value for a discrete set of storage levels.
\citet{alic2023optimal} apply SDP to a one-year problem with weekly
stages, uncertain inflows and prices, and environmental constraints that
depend on the reservoir level, and show that these constraints alter the
optimal production strategy. \citet{helseth2017iet} combine SDP with
weekly mixed-integer linear problems to retain non-convex operating
characteristics, and find that a linear version of their model
overestimates the profit from reserve sales. \citet{abgottspon2014} use
SDP between stages together with detailed subproblems within each stage.

The common strength of these applications is that SDP imposes no
convexity requirement on the stage problems, so detailed plant
descriptions can be kept while decisions still adapt to new information
every stage. The weakness is the discretisation of the state space: the
number of state combinations grows rapidly with the number of reservoirs
and other state variables. The reviewed SDP applications accordingly
concern individual plants or small cascades. The repeated solution of
stage problems for every storage level and uncertainty state also
becomes demanding when the stage problems themselves are detailed.

\subsubsection{Stochastic Dual Dynamic Programming}

Stochastic dual dynamic programming (SDDP), introduced by
\citet{pereira1991}, avoids discretising the state space by
approximating the future value function with cutting planes that are
refined iteratively. This allows multistage problems with several
reservoirs to be solved, and the method has become the standard approach
for large-scale hydropower scheduling under uncertainty
\citep{fullnerRebennack2025}. \citet{helseth2013} apply it to a
hydrothermal system that includes pumped storage and wind power, with
linearised transmission constraints, approximate start-up costs and a
fine time resolution, at the expense of a considerable computational
burden. They report that the approximated start-up costs and the
transmission constraints clearly affect the results.

The method relies on the structure of the problem. The convergence
analysis of \citet{shapiro2011} concerns multistage linear stochastic
programmes, and the cutting-plane argument does not extend directly to
non-convex stage problems. Price uncertainty requires particular care,
since uncertain prices enter the objective. \citet{helseth2016} handle
this by combining SDP for the price process with SDDP for the
reservoirs, which allows inflows, energy prices and reserve-capacity
prices to be uncertain. They also note that the required linear
relationships can introduce errors in the solution. This is the
counterpart of the SDP trade-off: SDDP scales to larger systems and
finer resolution, but restricts the physical description of the plant.

\paragraph{Extending the method}
Several studies seek to relax this restriction.
\citet{hjelmeland2019energies} include head-dependent, non-convex
generation when energy and reserve capacity are co-optimised, and
strengthen the cuts to approximate it. \citet{helseth2024approximating}
represent a detailed hydropower system within an aggregated model
through feasibility cuts computed before the SDDP run, building on
\citet{helseth2023hydropowera}. Both report better solutions together
with increased computation time. Such extensions narrow the gap between
the two dynamic-programming approaches without removing the underlying
trade-off, and none of them concerns pumped storage.

\paragraph{Applications to individual owners}
Applications of SDDP to an individual owner's valuation remain
concentrated on conventional reservoir hydropower.
\citet{bore2026assessing} use SDDP with a Markov chain for inflow and
prices to value generating-capacity upgrades across several markets.
For pumped storage, \citet{abgottspon2016multihorizon} use SDDP in a
multi-horizon model that links seasonal reservoirs with small daily
reservoirs, but in a deterministic case, and
\citet{schuller2024practical} use water values from an SDDP-based model
to value a pumping investment in the energy market. Owner-oriented
pumped-storage studies that include reserve markets have relied on SDP
or deterministic optimisation.

\paragraph{Physical detail and the choice of method}
The appropriate level of physical detail thus depends on the purpose of
the model. Short-term scheduling studies can afford discrete decisions
because their horizon is short, whereas studies of seasonal storage
simplify the plant to keep the problem tractable. The reviewed studies
document that simplifications affect the level of estimated income.
They do not examine whether simplifications also change the ranking of
alternative plant configurations, which is the relevant concern when
configurations are compared.

The methods thus form a progression in which each step gains in one
respect and gives up in another. Deterministic and scenario-based
formulations allow the most detailed plant description but do not yield
an operating policy under uncertainty. Rolling two-stage formulations
add uncertainty at the expense of its sequential resolution. SDP
represents sequential decisions with full physical detail for small
systems, and SDDP extends sequential decision-making to larger problems
and finer resolution under linearity requirements. The choice of method
consequently determines which of the features discussed in
Sections~\ref{sec:lit_valuation}--\ref{sec:lit_markets} a study is able
to represent.

% ---------------------------------------------------------------------
\subsection{Research Gap and Positioning of This Study}
\label{sec:lit_gap}

Table~\ref{tab:lit_comparison} summarises the reviewed studies along the
dimensions discussed above. The literature establishes several relevant
strands: economic valuation from different perspectives, alternative
representations of uncertainty and of the plant, participation in
multiple markets, comparison of plant configurations, and solution
methods with different capabilities. These strands address
complementary aspects of the value of pumped-storage flexibility.

The closest studies cover different parts of the capacity-valuation
question. \citet{abgottspon2014} connect pumped-storage scheduling with
ancillary-service opportunities, and \citet{helseth2016} connect
uncertain market conditions with water values. \citet{helseth2017iet}
additionally compare hydropower configurations, while
\citet{zisos2023mixing} investigate reservoir sizing. These
contributions provide a basis for analysing capacity interactions, but
their different assets, objectives and design choices do not establish
a common assessment of pump, turbine and reservoir capacities for a
pumped-storage owner.

The thesis by \citet{bore2026assessing} narrows this research need
further by directly connecting stochastic multi-market operation to the
value of generating-capacity upgrades. Consequently, the combination of
uncertainty, market participation and capacity valuation cannot itself
be treated as an unaddressed question. The remaining gap concerns the
joint contribution of pumping power, generating power and reservoir
capacity in a pumped-storage setting, together with the influence of
different price dynamics, which the reviewed studies have examined only
with known prices. The review also indicates that owner-oriented
pumped-storage studies with reserve markets have so far relied on
methods that either omit sequential uncertainty or are confined to small
state spaces.

Against this background, the present study asks how pumping power,
generating power and reservoir capacity, individually and in
combination, affect the owner's long-term operating value across
electricity markets, and how the relative value of these characteristics
changes with price volatility and persistence. The study adopts the
perspective of a price-taking owner and compares predefined plant
configurations under common market and resource conditions. In line
with the reviewed applications to reservoir hydropower, the problem is
solved with SDDP, using weekly stages and uncertain inflow and day-ahead
prices. The last row of Table~\ref{tab:lit_comparison} places the study
alongside the reviewed literature. The contribution sought is a
comparative understanding of capacity interactions under uncertainty; it
does not depend on claiming that stochastic scheduling, multi-market
participation or storage sizing is individually new.

% ---------------------------------------------------------------------
% Table 3. Check the row "This study" against ch. 6 when it changes
% (reserve products depend on Q-002).
\begin{landscape}
\begingroup
\footnotesize
\setlength{\tabcolsep}{4pt}
\renewcommand{\arraystretch}{1.15}
\begin{longtable}{@{}>{\raggedright\arraybackslash}p{0.12\linewidth}
                   >{\raggedright\arraybackslash}p{0.15\linewidth}
                   >{\raggedright\arraybackslash}p{0.19\linewidth}
                   >{\raggedright\arraybackslash}p{0.17\linewidth}
                   >{\raggedright\arraybackslash}p{0.18\linewidth}
                   >{\raggedright\arraybackslash}p{0.13\linewidth}@{}}
\caption{Reviewed studies compared by valuation perspective,
uncertainty, market participation, capacity question and solution
method, and the positioning of this study.}
\label{tab:lit_comparison}\\
\toprule
\textbf{Study} & \textbf{Perspective and asset} & \textbf{Time and uncertainty} & \textbf{Markets} & \textbf{Capacity question} & \textbf{Solution method}\\
\midrule
\endfirsthead
\multicolumn{6}{l}{\tablename~\thetable\ (continued)}\\
\toprule
\textbf{Study} & \textbf{Perspective and asset} & \textbf{Time and uncertainty} & \textbf{Markets} & \textbf{Capacity question} & \textbf{Solution method}\\
\midrule
\endhead
\midrule
\multicolumn{6}{r}{\footnotesize Continued on the next page}\\
\endfoot
\bottomrule
\addlinespace
\multicolumn{6}{@{}p{\linewidth}@{}}{\footnotesize SDP: stochastic
dynamic programming; SDDP: stochastic dual dynamic programming; LP:
linear programming; MILP: mixed-integer linear programming; NLP:
nonlinear programming; PHS: pumped-storage hydropower; DA/ID:
day-ahead/intraday; FCR: frequency containment reserve (N: normal
operation, D: disturbances); aFRR: automatic frequency restoration
reserve. Studies cited only for price behaviour or as general
methodological references are not included in the comparison.}\\
\endlastfoot
\citet{helseth2013} & System; hydrothermal system with pumped storage & Medium to long term; inflow, wind power and prices & System operation (no owner revenue) & Operation of a specified plant & SDDP\\
\citet{abgottspon2013stochastic} & Price-making owner; pumped storage & One year; inflow & DA and forward contracts & Operation of the studied plant & SDP\\
\citet{abgottspon2014} & Owner; pumped storage & One year; inflow and prices; known or uncertain week compared & Energy and ancillary services & Operation of the studied plant & SDP with stochastic or simplified within-week problems\\
\citet{helseth2016} & Price-taking owner; reservoir hydro & Medium term; inflow, energy and reserve-capacity prices & Sequential energy and reserve-capacity markets & Water values and schedules with reserve access & Combined SDP/SDDP\\
\citet{helseth2017iet} & Price-taking owner; reservoir hydro & One year; weekly Markov states & Day-ahead energy and primary reserve capacity & Alternative generator configurations & SDP with weekly MILP\\
\citet{chazarra2017} & Price-taking owner; closed-loop pumped storage & Short-term hourly scheduling; known inputs & Energy, secondary reserve and regulation energy & With and without hydraulic short-circuit operation & MILP\\
\citet{helseth2017} & System; Nordic power system & One year ahead each week; weather and exogenous prices & System operation (socio-economic surplus) & Transmission and pumped-storage expansions & Two-stage stochastic LP in a rolling horizon\\
\citet{ak2019quantifying} & Owner; cascade with pumping & Annual evaluation; historical inflows and price years & Energy revenues & Existing cascade versus pumped-storage operation & NLP with operating rules\\
\citet{dallinger2019socioeconomic} & System and owner; Austrian plant & One historical year; known inputs & DA and aFRR & Adding pump-turbines & Dispatch model and MILP\\
\citet{alic2023optimal} & Owner; pumped storage in a cascade & One year; weekly stages; inflow and energy prices & Price-based energy scheduling & Pumping addition and constraint effects & SDP\\
\citet{zisos2023mixing} & Hybrid island-energy design & Long-run simulation; wind, demand and conversion & Supply reliability and economic performance & Reservoir and PV sizing; fixed pump and turbine power & Simulation and optimisation\\
\citet{schuller2024practical} & Price-taking owner; Norwegian watercourse & Inflow and prices & Energy & A new plant with pumping & SDDP-based water values with hourly simulation\\
\citet{naval2025optimization} & Hybrid-system operator; pumped storage & Hourly operation over a year; wind & Grid exchange and renewable demand coverage & Operation of a specified hybrid configuration & MILP with chance constraints\\
\citet{bore2026assessing} & Price-taking owner; reservoir hydro & One year; weekly stages; inflow and prices; known within-week inputs & DA, ID, FCR-N, FCR-D and aFRR capacity; activation excluded & Generating-power upgrades; maximum investment cost & SDDP\\
\midrule
This study & Price-taking owner; pumped storage with natural inflow & Weekly stages; inflow and day-ahead prices; known within-week inputs & Day-ahead energy and reserve capacity; activation excluded & Combinations of pumping power, generating power and reservoir capacity & SDDP\\
\end{longtable}
\endgroup
\end{landscape}
